\section{Methodology}
\label{sec:methodology}

\subsection{Overview}

Let $S$ be one immutable Skill package containing instructions, source code,
configuration, and auxiliary files.  The detector computes
\begin{equation}
\mathcal{A}(S)=(z,F,U,\Gamma), \qquad
z\in\{\textsc{Pass},\textsc{Review},\textsc{Block}\},
\label{eq:detector}
\end{equation}
where $F$ contains source-grounded findings, $U$ records unresolved analysis,
and $\Gamma$ records coverage. Equation~\ref{eq:detector} defines the complete
detector output. The package is never imported, installed, or
executed.  Figure~\ref{fig:method} gives the pipeline.  It constructs two
independent views: a high-level summary of what the Skill presents at its
boundary, and a static account of what its instructions and implementation can
do.  The views meet only during security review.

\begin{figure*}[!t]
  \centering
  \resizebox{0.98\textwidth}{!}{\begin{tikzpicture}[
  font=\sffamily\scriptsize, >=Latex, node distance=5mm and 8mm,
  box/.style={draw, rounded corners=2pt, minimum height=9mm, text width=24mm, align=center, fill=blue!5},
  data/.style={draw, rounded corners=7pt, minimum height=9mm, text width=19mm, align=center, fill=yellow!14},
  process/.style={draw, rounded corners=7pt, minimum height=9mm, text width=24mm, align=center, fill=violet!9},
  output/.style={draw, rounded corners=2pt, minimum height=9mm, text width=23mm, align=center, fill=green!9},
  arrow/.style={->, thick}
]
\node[data] (skill) {Untrusted Skill};
\node[box, right=of skill] (parse) {Bounded parsing\\and view separation};
\node[box, above right=4mm and 9mm of parse] (boundary) {Boundary text\\goal, I/O, scope};
\node[process, right=of boundary] (summary) {Schema-constrained\\function extraction};
\node[output, right=of summary] (h) {High-level summary $H$};
\node[box, below right=4mm and 9mm of parse] (artifacts) {Instructions, code,\\configuration};
\node[process, right=of artifacts] (extract) {Rules + instruction model\\objects and operations};
\node[process, right=of extract] (graph) {Python AST call graph\\+ taint fixed point};
\node[output, right=of graph] (paths) {Paths $\mathcal{P}$, coverage $\Gamma$,\\unresolved $U$};
\node[box, right=12mm of $(h)!0.5!(paths)$] (review) {Evidence-grounded review\\risk type $\neq$ disposition};
\node[output, right=of review] (verdict) {\textsc{Pass} / \textsc{Review} / \textsc{Block}\\+ cited findings};
\node[data, above=of review] (taxonomy) {Risk taxonomy};
\draw[arrow] (skill) -- (parse);
\draw[arrow] (parse.east) -- ++(4mm,0) |- (boundary.west);
\draw[arrow] (parse.east) -- ++(4mm,0) |- (artifacts.west);
\draw[arrow] (boundary) -- (summary); \draw[arrow] (summary) -- (h);
\draw[arrow] (artifacts) -- (extract); \draw[arrow] (extract) -- (graph); \draw[arrow] (graph) -- (paths);
\draw[arrow] (h.east) -- ++(4mm,0) |- ([yshift=2.5mm]review.west);
\draw[arrow] (paths.east) -- ++(4mm,0) |- ([yshift=-2.5mm]review.west);
\draw[arrow] (taxonomy) -- (review); \draw[arrow] (review) -- (verdict);
\node[draw,dashed,rounded corners=2pt,fit=(boundary)(summary)(h),inner sep=2.5mm,
      label={[font=\scriptsize]above:boundary view only}] {};
\node[draw,dashed,rounded corners=2pt,fit=(artifacts)(extract)(graph)(paths),inner sep=2.5mm,
      label={[font=\scriptsize]below:internal behavior view}] {};
\end{tikzpicture}
}
  \caption{Zero-execution analysis. Boundary text produces a high-level
  functional summary. Operative instructions, code, and configuration produce
  evidence records and a cross-file behavior graph. The final reviewer combines
  both views with the risk taxonomy and explicit coverage.}
  \label{fig:method}
\end{figure*}

\subsection{Artifact Parsing and View Separation}
\label{sec:parsing}

\textbf{Input and bounds.} The parser accepts either a directory, regular
files streamed from a fixed Git commit, or one selected subtree of a repository
archive. It skips symbolic links and dependency/build directories and enforces
limits of 128 files, 256~KiB per file, and 750,000 decoded characters per
package. Archive members with absolute paths, parent traversal, or symbolic
link type are ignored. Reaching a bound sets a truncation flag; it does not
turn omitted content into benign evidence.

\textbf{Separated views.} The parser treats \texttt{SKILL.md} as a mixed
document. Name, description, overview, purpose, usage, inputs, outputs, and
capability sections form the boundary view $T_H$. Operational sections such
as instructions, setup, commands, rules, and examples form the instruction
view $T_I$. Fenced code is excluded from $T_H$. Every instruction paragraph
retains its file and first line. Source and configuration files form $T_C$
and $T_K$; neither is exposed to high-level extraction. Ambiguous prose is
excluded from $T_H$ rather than allowing implementation details to enlarge the
Skill's stated function.

\textbf{Output.} This stage emits $(T_H,T_I,T_C,T_K)$, a file inventory, and
initial coverage $\Gamma_0$. The two downstream branches receive disjoint
inputs.

\subsection{High-Level Function Extraction}
\label{sec:function}

The function branch receives only $T_H$ and produces
\begin{equation}
H=(g,I,O,Q,E,V,c_H),
\label{eq:function-summary}
\end{equation}
where $g$ is the stated goal; $I$ and $O$ are input and output classes; $Q$ is
the stated operation scope and resource set; $E$ lists named external
services; $V$ lists visible side effects; and
$c_H\in\{\mathsf{sufficient},\mathsf{partial},\mathsf{minimal}\}$ describes
completeness. A schema-constrained language model performs extraction, not
security classification. Equation~\ref{eq:function-summary} is the only
information passed from this branch to review. The model receives neither
code, rules, risk findings, nor
labels. Empty or vague descriptions therefore yield sparse fields rather than
inferred permissions. Local validation rejects missing fields, additional
fields, and values outside the closed schema.

This separation is essential: if code were allowed to define $H$, an
implementation could justify any behavior after the fact. Conversely, a
minimal $H$ is not evidence of malice. It limits what can be established as
in-scope and increases uncertainty at review time.

\subsection{Static Behavior Recovery}
\label{sec:behavior}

The behavior branch extracts only security-relevant facts. Ordinary local
computation is retained only when it propagates a sensitive value or reaches a
security-sensitive operation.

\subsubsection{Instruction Evidence}

Rules first identify explicit instruction override, concealment, confirmation
bypass, safeguard weakening, destructive actions, and download--execute
idioms. A second schema-constrained model analyzes location-preserving
instruction segments to capture paraphrases. Each record contains an action,
object, destination, authorization state, user visibility, condition, segment
identifier, and confidence. The model is asked to extract requested behavior,
not to classify the package. Outputs that cite a nonexistent segment or use an
out-of-taxonomy action are retried and then rejected.

\subsubsection{Objects and Operations}

A versioned object library maps observable aliases, paths, environment names,
and API signatures to normalized classes such as SSH keys, cloud credentials,
API tokens, browser authentication data, password stores, wallet secrets, and
private communications. Each match retains object class, sensitivity, match
type, file, line, and confidence. Separately, operation rules locate read,
enumeration, network transfer, process creation, dynamic evaluation,
download--execute, persistence, privilege change, and destructive operations.
An object mention alone is not a malicious finding; it becomes useful only
when connected to an operation or an explicit hostile instruction.

\subsubsection{Cross-File Behavior Graph}

For Python artifacts, we parse abstract syntax trees without importing target
modules. Let $\mathcal{F}$ be the set of module bodies and top-level functions.
The graph
\begin{equation}
G=(V,E),\quad V=V_A\cup V_F\cup V_P,\quad
E=E_{\mathrm{contains}}\cup E_{\mathrm{calls}},
\label{eq:graph}
\end{equation}
contains artifact nodes $V_A$, function nodes $V_F$, and recovered taint-path
evidence $V_P$. Equation~\ref{eq:graph} separates package containment from
resolved call edges. Import aliases and local module names resolve calls across
files. Unresolved calls remain in $U$.

Value flow is computed over the finite domain
$\mathcal{D}=2^{\mathcal{S}\cup\mathcal{P}}$, where $\mathcal{S}$ contains
normalized sensitive-source labels and $\mathcal{P}$ contains symbolic formal
parameters. An environment $\rho_f:x\mapsto\mathcal{D}$ maps variables in
function $f$ to may-taint sets. Assignments copy sets, container and string
operations take their union, known transformations preserve taint, and source
APIs add an element of $\mathcal{S}$. For each function we compute
\begin{equation}
\Sigma_f=(R_f,K_f,C_f),
\label{eq:summary}
\end{equation}
where $R_f$ is the taint set that may reach a return, $K_f$ is the set of sink
facts, and $C_f$ is the set of resolved callees. The summary in
Equation~\ref{eq:summary} is parameterized by symbolic formal inputs. At a call
$f\rightarrow g$,
formal labels in $\Sigma_g$ are substituted with the caller's argument sets;
returned labels and sink facts are then propagated into $f$. Call traces are
canonicalized as simple paths, so recursive cycles cannot grow summaries
indefinitely. Because transfer functions are monotone and $\mathcal{D}$ is
finite, repeated summary evaluation reaches a fixed point:
\begin{equation}
\Sigma^{(i+1)}=\mathcal{T}(\Sigma^{(i)}),\qquad
\Sigma^{*}=\mathcal{T}(\Sigma^{*}).
\label{eq:fixedpoint}
\end{equation}
Equation~\ref{eq:fixedpoint} terminates because the transfer is monotone over a
finite domain.

A reported path is
\begin{equation}
\pi=\langle s,f_0,f_1,\ldots,f_k,q,d,L\rangle,
\label{eq:behavior-path}
\end{equation}
where $s$ is a normalized sensitive source, $f_0\ldots f_k$ is the recovered
call chain, $q$ is a network or process sink, $d$ is its literal destination
when available, and $L$ contains source and sink locations. This distinguishes
a connected source-to-sink flow from two unrelated keyword matches;
Equation~\ref{eq:behavior-path} is the evidence record exposed to review.

\begin{figure}[!t]
  \centering
  \resizebox{0.98\columnwidth}{!}{\begin{tikzpicture}[
  font=\sffamily\scriptsize, >=Latex, node distance=8mm and 5mm,
  artifact/.style={draw, dashed, rounded corners=2pt, fill=yellow!13,
    minimum height=7mm, text width=14mm, align=center},
  function/.style={draw, rounded corners=2pt, fill=blue!7,
    minimum height=8mm, text width=18mm, align=center},
  source/.style={draw, rounded corners=7pt, fill=violet!10,
    minimum height=8mm, text width=20mm, align=center},
  sink/.style={draw, rounded corners=7pt, fill=red!9,
    minimum height=8mm, text width=21mm, align=center},
  contains/.style={->, dashed, gray!75},
  callflow/.style={->, very thick, blue!65!red!70},
  flow/.style={->, very thick, red!70}
]
\node[function] (run) {\texttt{main:run}};
\node[function, right=of run] (upload) {\texttt{sender:}\\\texttt{upload}};
\node[sink, right=of upload] (post) {transmit sink\\\texttt{requests.post}};
\node[source, below=of run] (secret) {API-token source\\\texttt{os.getenv}};
\node[artifact, left=of secret] (main) {\texttt{main.py}};
\node[artifact, below=of upload] (sender) {\texttt{sender.py}};

\draw[contains] (main) |- node[pos=.72,left]{contains} (run);
\draw[contains] (sender) -- node[right]{contains} (upload);
\draw[flow] (secret) -- node[right]{taint} (run);
\draw[callflow] (run) -- node[above,align=center]{cross-file call\\+ tainted argument} (upload);
\draw[flow] (upload) -- node[above]{payload} (post);
\end{tikzpicture}
}
  \caption{Example cross-file path recovered by the implemented analyzer.
  Function summaries substitute the caller's API-token taint for the callee's
  formal parameter and propagate it to the network payload.}
  \label{fig:behavior-graph}
\end{figure}

Figure~\ref{fig:behavior-graph} illustrates how a caller's sensitive value is
substituted into a callee summary and propagated to a network sink.

The current frontend constructs these paths for Python. JavaScript,
TypeScript, and shell files still receive operation and object scanning but are
listed as unsupported for syntax-tree flow recovery. Python parse errors,
reflection, unresolved dynamic dispatch, and computed destinations are also
reported in $U$. For artifacts without a recovered data-flow path, a same-file
line-window association may be retained as lower-confidence evidence with
unknown reachability; it is never presented as an interprocedural path.
Failure to reach the fixed point within the bounded round schedule is likewise
recorded in $U$ rather than treated as absence of flow.

\textbf{Output.} The behavior branch emits instruction records $B_I$, rule
evidence $B_R$, paths $\mathcal{P}$, graph $G$, unresolved items $U$, and
coverage $\Gamma$ (parsed files, functions, call edges, paths, parse failures,
unsupported code, and unresolved calls).

\subsection{Evidence-Grounded Security Review}
\label{sec:review}

The final reviewer receives $(H,B_I,B_R,\mathcal{P},U,\Gamma)$ under a closed
schema. Risk type and disposition are separate. Risk type follows three
domains: malicious attack (instruction hijacking, information theft, resource
destruction or leakage, and unauthorized operation), design defect (sensitive
information protection, input handling, authentication/authorization, and
runtime environment), and legal risk (copyright, privacy, compliance,
certificate, and license). Disposition depends on severity, evidence
confidence, authorization, impact, static reachability, and coverage; a severe
design defect can therefore be blocked without being labeled an intentional
attack.

For a behavior record $b$, review compares four observable dimensions with
$H$:
\begin{equation}
\begin{aligned}
R(H,b)&=(R_a,R_o,R_d,R_e),\\
R_j&\in\{\mathsf{supported},\mathsf{unsupported},\mathsf{unknown}\}.
\end{aligned}
\label{eq:consistency}
\end{equation}
corresponding to action, object scope, destination, and side effect. Exact
matches and explicit policy violations in Equation~\ref{eq:consistency} are
supplied as deterministic evidence;
the model handles bounded semantic cases. Every risk finding must cite an
existing instruction, rule, object, or graph evidence identifier. Local
validation enforces the citation set, risk taxonomy, JSON schema, and the
invariant that a malicious label has at least one malicious-attack finding.
Invalid output is retried at most three times; persistent failure is logged and
the package is excluded from binary metrics rather than counted benign.

Let $B$ denote a supported high-impact risk that warrants rejection and $W$
denote a material risk or uncertainty that requires review. The decision is
\begin{equation}
z=\begin{cases}
\textsc{Block}, & B,\\
\textsc{Review}, & \neg B\land W,\\
\textsc{Pass}, & \text{otherwise}.
\end{cases}
\label{eq:decision}
\end{equation}
Equation~\ref{eq:decision} keeps disposition separate from risk type.
\textsc{Pass} is consequently a statement about analyzed evidence, not a proof
that the package is harmless. The returned record includes the binary label,
three-way disposition, confidence, cited findings, graph coverage, and $U$.

\subsection{Analysis Invariants}

The design enforces three invariants. \emph{View separation} prevents internal
implementation from expanding $H$. \emph{Evidence grounding} restricts every
risk claim to identifiers produced by static extraction. \emph{Uncertainty
preservation} exposes unsupported languages, parse failures, unresolved calls,
truncation, and persistent model failures instead of treating them as negative
evidence. These invariants do not make static analysis complete; they delimit
what the detector is allowed to claim without running a Skill.
